\backhead{FURTHER EXERCISES}{khred}
\begin{exlist}
\item A project uses the loam with a 300 cm trigger. Give the sequence of layers you would open to answer: (a) how much water is saved by a mulch; (b) where the saved water goes; (c) whether the saving reduces irrigation.\exans{(a) Mulch against bare soil; (b) Fate of saved water; (c) Field balance with and without mulch}
\item The Dashboard reports a yield of 6.1 Mg\,ha$^{-1}$, irrigation 140 mm and rainfed yield 4.3 Mg\,ha$^{-1}$. Find the irrigation water use efficiency in kg\,m$^{-3}$.\exans{1.29 kg\,m$^{-3}$}
\item A layer shows 20 years of results. If each year has 365 daily values for 8 variables, how many values does the layer draw? Why does the Dashboard show annual means instead?\exans{58\,400; to keep the reading clear}
\item Two layers give the drainage of the same soil as 160 and 175 mm. List three reasons that may explain the difference.\exans{different engines, inputs or periods}
\item Explain, with numbers, how the change of soil from the loam ($\TAW=150$ mm) to the sand ($\TAW=60$ mm) changes the interval between irrigations at $ET_c=5$ mm\,d$^{-1}$ and $p=0.5$.\exans{15 d to 6 d}
\end{exlist}
\backhead{ESSAY QUESTIONS}{khred}
{\small\itshape Answer each question in about one page, with equations, a sketch and at least one number from this chapter.}\par\vspace{3pt}
\begin{exlist}
\item Describe the organisation of SWAFL into layers. Why is a layered organisation useful for teaching and for decision making?
\item Explain how the layers share one project and how a change in one layer appears in another.
\item Compare the answers of the layers with those of a new Daisy run. When is each the right tool?
\item Discuss the role of the Dashboard as the common output of all layers.
\item Why does each number on the screen carry its equation? Discuss the value of equation traceability for students and for advisers.
\item Choose one layer and describe how you would use it to advise a farmer, step by step.
\end{exlist}
\dykband{EXERCISES IN SWAFL}
{\small Open SWAFL (desktop or online) with the benchmark project. The chapter of the \href{https://github.com/halilunaziru73-creator/Coupled-Water-Fluxes-Pore-Scale-Dynamics-Precision-Irrigation-Mulching-Feedback-Loops/raw/main/Software/v1.1.0/SWAFL_Desktop/docs/manual_v1.1.0/SWAFL_Manual_v1.1.0.pdf}{SWAFL manual} that explains each task is given in brackets; record your values in a table and compare them with what is expected (\emph{Check}).}\par\vspace{3pt}
\begin{exlist}
\item {\color{khblue}[Manual, Chapter 2]} Open SWAFL and the welcome dialog; start a new project and open each group of the Layers panel in turn. List the layers of each group.\par\noindent{\small\itshape Check: the groups and layers listed in the manual.}
\item {\color{khblue}[Manual, Chapter 3]} Change three Properties (soil, crop, trigger) and observe which layers redraw.\par\noindent{\small\itshape Check: every layer that depends on the property redraws.}
\item {\color{khblue}[Manual, Chapter 14]} Open the four tabs of the Dashboard and describe what each shows for the field balance.\par\noindent{\small\itshape Check: results, design, charts and messages.}
\item {\color{khblue}[Manual, Chapter 16]} Export the results of one layer and open them in a spreadsheet.\par\noindent{\small\itshape Check: the exported file holds the values shown on screen.}
\end{exlist}
